%------------------------------------------------
% Latex-Grundgerüst für Seminarausarbeitungen
%
% zu Erzeugen mit
%   pdflatex ausarbeitung.tex
%
% von Carsten Gutwenger, Nils Kriege
% 
%------------------------------------------------

\documentclass[a4paper,12pt,twoside]{article}
\usepackage[a4paper,left=3.5cm,right=2.5cm,bottom=3.5cm,top=3cm]{geometry}

% Encoding der .tex-Datei
% muss je nach Editor-Einstellungen eventuell angepasst werden!
%   latin1    Unix/Windows
%   ansinew   Alternative für Windows
%   applemac  Apple
%   utf8      Alternative für Unix/Windows/Apple
\usepackage[utf8]{inputenc}

\usepackage[ngerman]{babel}                                 % deutsche Sprache
\usepackage[T1]{fontenc}                                    % Unterstützung für Umlaute mit Fonts
\usepackage[pdftex]{graphicx,color}                         % Einbetten von Grafiken, Farbe
\usepackage[format=plain,small,bf,indention=7.5mm]{caption} % verbesserte Beschriftung von Abbildungen
\usepackage{subfig}                                         % subfigures
\usepackage{url}                                            % \url{}-Kommando
\usepackage{fancyhdr,float}                                 % Kopf-/Fußzeilen
\usepackage[pdftex,pdfpagelabels]{hyperref}                 % Hyperlinks im PDF
\usepackage[round]{natbib}                                  % Literaturverzeichnis
\usepackage{doi}                                            % DOIs
\usepackage[section,boxed]{algorithm}                       % Floating-Umgebung für Algorithmen
\usepackage{algpseudocode}                                  % Pseudocode für Algorithmen
\usepackage[standard]{ntheorem}                             % Theorem-Umgebung
\usepackage{mathtools}										% Tools zum Anzeigen von Formeln

\captionsetup{farskip=10pt,topadjust=0pt,captionskip=10pt,nearskip=0pt,margin=10pt}

\pagestyle{fancy}
\fancyhead[LE,RO]{\thepage}
\fancyhead[RE]{\nouppercase{\slshape \leftmark}}
\fancyhead[LO]{\nouppercase{\slshape \rightmark}}
\fancyfoot[C]{}


\begin{document}

%------------------------------------------------
% Titelseite
%------------------------------------------------

\begin{titlepage}
\vspace*{-2cm}
\newlength{\links}
\setlength{\links}{-1.5cm} \sf \LARGE

\hspace*{\links}
\begin{minipage}{12.5cm}
\includegraphics[width=8cm]{tud_logo_rgb}
\end{minipage}

\vspace*{4cm}

\large
\begin{center}
{\Large Seminarausarbeitung} \\[1ex]
{\LARGE\textbf{On the Power of Color Refinement}}\\[3ex]
Florian Lüdiger\\[1ex]
\today\\[7ex]
im Rahmen des Seminars\\[1ex]
{\Large\textbf{Algorithm Engineering}}\\[1ex]
von Prof.~Dr.~Petra Mutzel\\[1ex]
Wintersemester 2017/18
\end{center}

\vspace*{4cm}
\hspace*{\links}
\begin{minipage}[b]{15cm}
\normalsize \raggedright
\textbf{Betreuer:} \\
Christopher Morris \\[2ex]

\textbf{Basierend auf:}\\
V. Arvind, Johannes Köbler, Gaurav Rattan und Oleg Verbitsky\\
On the Power of Color Refinement\\
\url{https://link.springer.com/chapter/10.1007/978-3-319-22177-9_26}
\end{minipage}

\definecolor{TUGreen}{rgb}{0.517,0.721,0.094}
\vfill
\hspace*{\links}
\begin{minipage}[b]{8cm}
\normalsize \raggedright
Fakultät für Informatik\\
Algorithm Engineering (Ls11)\\
Technische Universität Dortmund
\end{minipage}

\end{titlepage}

%------------------------------------------------
% Inhaltsverzeichnis
%------------------------------------------------

\tableofcontents
\clearpage


%------------------------------------------------
% Einführung
%------------------------------------------------

\section{Einführung}
\label{sec/einfuehrung}

\begin{itemize}
	\item Allgemeine Informationen zum Paper
	\item Was ist das Ziel der Forschung und inwiefern wurde dieses erreicht?
	\item Überblick über die Kapitelstruktur
\end{itemize}


\section{Graph-Isomorphie und Color Refinement}


\section{Begriffserklärung}

Einführung, Wiederholung und Erklärung folgender Begriffe und Notationen:
\begin{itemize}
	\item zugänglicher Graph
	\item Unigraph
	\item Bipartites Komplement
	\item Biregulärer Graph
	\item Knotenfärbung
	\item Partition
	\item gerechte Partition
	\item Matching Graph
	\item Zellgraph
	\item Eigenschaften von Zellgraphen
\end{itemize}

Gegebenenfalls wird dieses Kapitel nicht explizit umgesetzt, sondern die Begriffe an der entsprechenden Stelle im Text eingeführt, an der sie gebraucht werden.

\section{Voraussetzungen für zugängliche Graphen}

\begin{itemize}
	\item Erklärung der lokalen Struktur zugänglicher Graphen (A,B)
	\item Erklärung der globalen Struktur zugänglicher Graphen (C,D,E,F,G,H)
	\item jeweils gegebenenfalls mit ausführlicher Erklärung der Bedeutung und Beispielen
\end{itemize}

\section{Korrektheitsbeweis}

\begin{itemize}
	\item Fokussierung auf die Beweisidee für die jeweiligen Lemmata
	\item Ausführlichere und leichter verständliche Aufbereitung der Beweise aus dem Paper
\end{itemize}

\section{Laufzeitanalyse}
\label{sec/laufzeit}

\begin{itemize}
	\item Beweis der Laufzeit von $O((n+m)\log n)$
\end{itemize}

\section{Fazit}

\begin{itemize}
	\item Zusammenfassung der Ergebnisse
	\item Bezugnehmen auf die Einleitung und ob die Ziele erreicht wurden
	\item Abschluss der Arbeit
\end{itemize}


%\subsection{Beispiele für Referenzen}
%
%\begin{itemize}
% \item Die starken Zusammenhangskomponenten eines gerichteten Graphen können in Linearzeit bestimmt werden~\citep{Tar72}.
% \item \citet{Tar72} hat gezeigt, dass \dots
%\end{itemize}
%
%
%\subsection{Theorem-Umgebung}
%
%% Vordefinierte Umgebungen:
%% Theorem, Lemma, Proposition, Satz, Korollar, Definition, Beispiel, Anmerkung, Bemerkung, Beweis
%\begin{Theorem}[Optionaler Titel]
% Aussage.
%\end{Theorem}


%------------------------------------------------
% Literaturverzeichnis
%------------------------------------------------

\bibliographystyle{abbrvnat-ger}
\bibliography{literatur}
\addcontentsline{toc}{section}{\bibname}


\end{document}
